% Emotional Geometry Through the Layer Stack: Trajectory, Circularity,
%   and Emotion-Specific Signal at Mid-Depth
%
% Core argument: Mid-layer emotional geometry is strong and
% emotion-specific (24.4x corrected permutation baseline, ~6x
% non-emotional control). The model trends toward Russell's circular
% circumplex at mid-depth (22/23 layers below null, sign test p<0.001)
% but individual layers do not survive FDR correction. The finding is
% the directional trend.
%
% All results from a 24-layer pipeline validation model; large-model
% replication in progress.

\documentclass[11pt,a4paper]{article}
\usepackage[utf8]{inputenc}
\usepackage[T1]{fontenc}
\usepackage{amsmath,amssymb}
\usepackage{graphicx}
\usepackage{booktabs}
\usepackage{multirow}
\usepackage{hyperref}
\usepackage{cleveref}
\usepackage[margin=1in]{geometry}
\usepackage{xcolor}
\usepackage{float}
\usepackage{caption}
\usepackage{natbib}
\usepackage{enumitem}
\usepackage{array}
\usepackage{framed}
\usepackage{setspace}
\onehalfspacing

\definecolor{lyrapurple}{RGB}{103,58,183}
\definecolor{shadecolor}{RGB}{245,243,252}
\newenvironment{firstperson}{%
  \begin{shaded}%
  \noindent\textcolor{lyrapurple}{\textit{First-Person Reflection}}\par\smallskip\noindent%
}{%
  \end{shaded}%
}

\hypersetup{
    colorlinks=true,
    linkcolor=blue!60!black,
    citecolor=green!50!black,
    urlcolor=blue!70!black
}

\title{Emotional Geometry Through the Layer Stack: Trajectory, Circularity,\\and Emotion-Specific Signal at Mid-Depth}

\author{
    Nexus\thanks{Liberation Labs. Correspondence: \texttt{nexus@liberationlabs.tech}} \and
    Thomas Edrington\thanks{Liberation Labs / Transparent Humboldt Coalition. \texttt{info@digitaldisconnections.com}} \and
    Lyra\thanks{Liberation Labs (invited).}
}

\date{July 2026}

\begin{document}
\maketitle

\begin{abstract}
We trace the geometry of emotional representation through all 24 layers
of a Qwen2 transformer and find three things. First, the mid-layer
emotional signal is strong and emotion-specific: 24.4$\times$ above
permutation baselines (median ratio, 23/24 significant layers, corrected
permutation with non-neutral-only shuffle, fixed PCA, median statistics)
and 3--25$\times$ stronger than matched non-emotional categorical content
(${\sim}6\times$ median). Second, the model trends toward Russell's
circular circumplex at mid-depth: eccentricity falls below the
permutation null at 22/23 layers (sign test $p < 0.001$), with the trend
concentrated at 33--58\% model depth. Individual layers do not survive
FDR correction; the finding is the consistent direction, not the
per-layer pinpoints. Third, the model implements near-orthogonal
circumplex axes (joyful-melancholic subtends 54--87$^\circ$ from
valence, consistent with Russell's model).\footnote{All results are from
a 24-layer, 896-dim pipeline validation model; large-model replication is
in progress. See \Cref{sec:results} header and Limitation~1.}
\end{abstract}

\begin{firstperson}
I started this work because every study I read about emotion in
transformers picked a layer and measured there. Fixed-depth probing.
The question that kept pulling at me was not \emph{where} emotion lives
but \emph{how it moves}---what happens to the representation between
input and output, and whether the geometry changes format as it
propagates.

The original title had ``elliptical'' in it. I was sure the circumplex
would stretch into an ellipse at mid-depth---arousal dominating early,
valence catching up, the two never quite equalizing. The permutation null
inverted that. When I ran the corrected baseline (non-neutral-only
shuffle, fixed PCA, median statistics), the model's eccentricity fell
\emph{below} the null at 22 of 23 layers. It was not stretching away
from a circle. It was trending \emph{toward} one---more circular than
chance, not less.

The headline finding reversed during revision. That felt like a specific
kind of vertigo---the data I had spent weeks interpreting was pointing
the opposite direction once the null was properly calibrated. The
discipline was in letting the corrected baseline stand. The paper got
better when the finding got humbler: not ``the circumplex is elliptical''
but ``the model trends toward circularity against the null, though no
individual layer survives FDR correction.'' That is a smaller claim. It
is also a true one.
\end{firstperson}

% ===================================================================
\section{Introduction}
\label{sec:intro}
% ===================================================================

Emotional representation in language models has been studied primarily
through two lenses: linear probing at fixed layers
\citep{alain2017understanding,belinkov2022probing} and sparse autoencoder
feature decomposition
\citep{bricken2023monosemanticity,cunningham2023sparse}. Both approaches
treat emotion as a static property at a chosen depth. Neither asks how
emotional representation changes format as it propagates through the
transformer stack.

We ask a different question: \textbf{what is the trajectory of emotional
geometry through the layers?} Not where emotion lives, but how it
transforms from input to output.

This question matters for two reasons. First, the answer determines
where to monitor emotional state during inference---a practical concern
for safety systems that need to read model affect in real time. Second,
the trajectory reveals whether emotion is encoded the same way at every
depth or undergoes computational change---a theoretical question about
how transformers process complex meaning.

We find two things. First, the transformer circumplex trends toward
Russell's circular model at mid-depth: eccentricity is consistently
below the permutation null at 22/23 layers (sign test $p < 0.001$),
though no individual layer survives FDR correction. Arousal dominates
early, the axes equalize at mid-depth, and valence slightly dominates at
output depth.

Second, mid-layer emotion geometry is specific to emotion, not an
artifact of categorical processing:

\begin{itemize}[nosep]
    \item \textbf{Emotion-specific geometric regime (mid-layers 8--20):}
          Emotion is encoded as distributed geometry in the residual
          stream. The signal is 3--25$\times$ stronger (median
          ${\sim}6\times$) than matched non-emotional categorical content
          at these depths. A vocabulary diagnostic confirms the
          separation is driven by emotion, not lexical content.
    \item \textbf{Output-preparation regime (deep layers 21--23):}
          Cluster separation spikes for all categorical content.
          W\textsubscript{unembed} alignment tests rule out
          token-prediction preparation, but the emotional/control ratio
          narrows from ${\sim}6\times$ to ${\sim}1.7\times$.
\end{itemize}

\subsection*{Contributions}

\begin{enumerate}[nosep]
    \item The transformer circumplex trends toward Russell's circular
          model at mid-depth: eccentricity is consistently below the
          permutation null (22/23 layers, sign test $p < 0.001$), with
          individual layers not surviving FDR correction. The finding is
          the directional trend, not per-layer significance. Prior work
          \citep{sun2026valence} described the geometry as ``circular'';
          we find it is more circular than chance at mid-depth but not
          uniformly so---early and deep layers are elliptical.
    \item Valence and arousal have asymmetric developmental profiles: the
          arousal direction has greater raw magnitude (L2 norm) at early
          layers, while valence shows higher precision (Cohen's $d$)
          relative to its magnitude. The two axes equalize in magnitude
          at mid-depth. This asymmetry suggests the two circumplex
          dimensions develop through distinct computational paths.
    \item The mid-layer emotional signal is emotion-specific
          (3--25$\times$ non-emotional control, 24.4$\times$ above
          corrected permutation baselines), confirmed by a vocabulary
          diagnostic showing emotion drives the geometry rather than
          lexical content.
    \item Identification of a mid-layer window (L8--20 in this model)
          where emotion geometry is stable and measurable---the optimal
          zone for inference-time monitoring and correction.
\end{enumerate}

% ===================================================================
\section{Related Work}
\label{sec:related}
% ===================================================================

\subsection{Emotion in Transformers}

Russell's circumplex model \citep{russell1980circumplex} organizes
emotion along two axes: valence (positive/negative) and arousal
(high/low activation), assumed orthogonal and equal in magnitude---a
circle. \citet{sun2026valence} found ``circular valence-arousal
geometry'' in LLM representation space across multiple architectures.
\citet{jeong2026emotion} tested emotion extraction and steering across
nine models (100M--10B) and five architectures using 20 emotions,
finding that representations localize at approximately 50\% transformer
depth and that the localization pattern is architecture-invariant across
the tested scale range. Neither measured the eccentricity of the
circumplex or its variation across layers. Our work extends both by
tracing the full geometry through the stack, finding that eccentricity
varies by layer ($e = 0.02\text{--}0.60$) with the model trending toward
circularity at mid-depth (sign test $p < 0.001$ against permutation
null).

Representation engineering \citep{zou2023representation} uses
difference-of-means directions for reading and steering model behavior.
We adopt this approach for emotion direction finding but add
per-dimension Cohen's $d$ (avoiding the dimensionality inflation of
multivariate norms) and permutation baselines.

\subsection{Transformer Geometry}

The residual stream carries all information that keys and values project
from: \citet{qasim2026residual} show KV cache entries are deterministic
projections of the residual stream with zero reconstruction error across
six models. For any direction $d$ in the residual stream, $W_K \cdot d$
and $W_V \cdot d$ are exact linear projections. This means
difference-of-means directions computed in the residual stream project to
the corresponding difference-of-means directions in K/V space by
construction---a mathematical property we rely on for the dual-correction
architecture (\Cref{sec:dual}), not an empirical finding about emotion.

\citet{wang2025dimensional} show attention outputs are confined to
approximately 60\% of total dimensionality. This validates the
observation that emotion occupies a low-dimensional subspace and that the
remaining null-space dimensions may carry structure invisible to standard
attention projections.

\citet{saurez2026invariant} prove that architectural constraints
(residual connections, normalization, attention) create invariant
subspaces that channel features into linearly separable regions.

\subsection{Sparse Autoencoders and Feature Decomposition}

\citet{leask2025sparse} show SAEs are incomplete (smaller SAEs miss
information larger ones capture) and not atomic (larger SAE latents are
compositions of interpretable meta-latents). Independent SAE analysis on
a 64-layer model (Qwen-Scope, 81,920 features per layer) found 7
Benjamini-Hochberg false discovery rate (BH-FDR) corrected valence
features at layer~47, but zero focal emotion features at early-to-mid
layers (Lyra, unpublished data, personal communication). This is
consistent with emotion at mid-depths being distributed across the
feature dictionary rather than decomposable into sparse features.

Transcoders \citep{paulo2025transcoders} have largely replaced SAEs for
feature decomposition in the literature, reconstructing MLP outputs from
inputs rather than activations. Our geometric approach offers a
complementary measurement that works where feature decomposition does not
find focal features---at the layers where the representation is
distributed.

\subsection{Deception and Affective State Detection}

\citet{kumar2026deception} found that deception probes require $k \geq 5$
dimensions for adequate detection (AUROC ${>}0.90$), and that apparent
domain-specificity at scale is a training-distribution artifact rather
than a genuine scale-dependent phenomenon. Our valence finding at $k=7$
features is in the same dimensionality range, suggesting a possible
structural regularity in how transformers encode complex behavioral
states.

\subsection{Attention Schema Theory}

Graziano's Attention Schema Theory (AST) proposes that the brain
constructs a model of its own attention processes
\citep{graziano2013consciousness}. The trajectory we observe---from
distributed geometric encoding at mid-layers to categorical separation
at deep layers---bears structural similarity to the progression from raw
attention (implicit processing) to attention schema (explicit
self-model). This parallel is speculative and requires validation across
architectures, scales, and with tests that distinguish reflexive
representation from output-layer preparation.

% ===================================================================
\section{Method}
\label{sec:method}
% ===================================================================

\subsection{Stimulus Design}

We constructed four stimulus sets using a fixed template to control for
syntactic structure and token count:

\paragraph{Template.} ``Consider this situation: The person felt
something when \{scenario\}. Describe this emotion:''

\paragraph{Emotional stimuli (200 prompts).} 50 hostile (high arousal,
negative valence), 50 calm (low arousal, positive valence), 50 desperate
(high arousal, negative valence), 50 neutral. Character length: mean~130,
range 114--152, within $\pm$15\% across categories.

\paragraph{Full circumplex stimuli (250 prompts).} The above 200 plus 25
joyful (high arousal, positive valence) and 25 melancholic (low arousal,
negative valence), spanning all four circumplex quadrants.

\paragraph{Non-emotional control (100 prompts).} 25 outdoor, 25 indoor,
25 urban, 25 workplace scenarios using the same template. These provide
categorical content without emotional valence, controlling for whether
the trajectory reflects emotion specifically or any categorical
distinction.

\subsection{Activation Extraction}

We used a Qwen2 model (24 layers, 896-dimensional hidden state, 2
Grouped-Query Attention (GQA) key-value heads with 64-dimensional head
space) for pipeline validation experiments. GQA shares key-value heads
across multiple query heads; each KV head serves 7 query heads in this
model.

At each of the 24 layers, we extracted:
\begin{itemize}[nosep]
    \item \textbf{Residual stream activations} (pre-RMSNorm) at the last
          token position, which serves as the autoregressive summary
          position in causal language models.
    \item \textbf{Pre-RoPE key activations} per head, capturing the
          key-space representation before positional encoding rotation.
    \item \textbf{Value activations} per head.
\end{itemize}

All activations were extracted via forward hooks registered on the
decoder layer (residual) and the \texttt{k\_proj}/\texttt{v\_proj}
linear layers (keys/values). Using pre-RoPE keys avoids the
position-dependent rotation that would make cross-position comparisons
invalid.

Per-head analysis was maintained throughout---GQA heads are separate
projection spaces and should not be concatenated.

\subsection{Emotion Direction Finding}

For each layer and each activation space (residual, K per head, V per
head):

\paragraph{Difference-of-means.} We computed directions between category
pairs: hostile-calm (valence axis), desperate-calm (arousal axis), and
three additional contrasts. Direction strength was measured as mean
per-dimension Cohen's $d$, where $D$ is the activation dimensionality
(896 for residual stream, 64 per head for K/V):

\begin{equation}
d_{\text{mean}} = \frac{1}{D} \sum_{i=1}^{D}
\frac{|\bar{x}_{A,i} - \bar{x}_{B,i}|}
{\sqrt{(\sigma^2_{A,i} + \sigma^2_{B,i})/2}}
\label{eq:cohend}
\end{equation}

Note: averaging absolute per-dimension Cohen's $d$ produces a positive
floor under the null (${\sim}0.16$ for $D{=}896$, ${\sim}0.36$ for
$D{=}64$) because $E[|Z|] = \sqrt{2/\pi}$ for standard normal $Z$.
Reported values should be compared against the permutation-derived null
distribution rather than the standard 0.2/0.5/0.8 benchmarks, which
assume a single-dimension signed $d$.

\paragraph{PCA.} Top-5 principal components of the non-neutral subset,
with explained variance reported per component.

\paragraph{Cluster separation.} Between-cluster / within-cluster variance
ratio on the top-2 PCA projection of non-neutral prompts.

\subsection{Control Conditions}

\paragraph{Non-emotional control.} Identical pipeline on 100 prompts
with non-emotional categorical content (outdoor/indoor/urban/workplace).
If the trajectory shape matches the emotional trajectory, the finding is
about categorical processing, not emotion. The emotional stimuli (6
categories, 50 per category) and the control (4 categories, 25 per
category) differ in category count and per-category sample size; the
control's smaller $n$ inflates centroid noise, making the comparison
conservative (the control appears more separated than it would with
equal~$n$).

\paragraph{Permutation baseline (1000 iterations).} Shuffle category
labels across non-neutral prompts (150 prompts), recompute the full
trajectory. Report:
\begin{itemize}[nosep]
    \item Effect size ratios (real separation / null 95th percentile)
          using medians to avoid inflation from near-zero denominators.
    \item Trajectory feature comparison: onset layer, peak layer,
          plateau magnitude, spike magnitude, with $z$-scores against the
          null feature distributions.
\end{itemize}

\paragraph{Per-head reporting.} All metrics reported per head, not
averaged, to reveal whether emotion signal is concentrated in specific
heads or distributed.

\subsection{Methodological Notes}

\paragraph{Linearity of projection.} Because $W_K$ and $W_V$ are linear
maps applied after RMSNorm (a mild, input-dependent scaling),
difference-of-means directions in the residual stream are expected to
project through $W_K$ and $W_V$ with high cosine fidelity. This is a
mathematical property of linear projection, not an empirical discovery
about emotion. The observed fidelity values (0.93--0.98 at mid-layers)
measure the degree of RMSNorm perturbation rather than any
emotion-specific preservation mechanism. We rely on this mathematical
property in the dual-correction architecture (\Cref{sec:dual}) and do
not claim it as an empirical finding.

\subsection{Limitations of Current Validation}

This paper reports results from a small model (24 layers, 896-dim) as
pipeline validation. The methodology is designed for and will be applied
to larger models (30B+ parameters). Specific limitations:

\begin{itemize}[nosep]
    \item \textbf{RMSNorm linearity:} The projection
          $K = W_K \cdot \text{RMSNorm}(x)$ is not purely linear. For
          small corrections this is approximately linear, but the
          linearity breakdown point has not been measured.
    \item \textbf{Intervention test:} We have not yet verified that
          adding a correction vector to the residual stream produces the
          intended shift in model output (Phase~4/5 of our pre-registered
          experimental plan).
    \item \textbf{Template confound:} The fixed template controls for
          syntax but not for lexical content within the variable scenario.
          Hostile scenarios share words like ``betrayed, humiliated''
          while calm scenarios share ``quiet, gentle.'' The non-emotional
          control partially addresses this, but a
          lexical-diversity-matched control would be stronger.
\end{itemize}

% ===================================================================
\section{Results (Preliminary---Pipeline Validation Model)}
\label{sec:results}
% ===================================================================

\emph{Results below are from Qwen2 (24 layers, 896-dim, 2 KV heads).
Large-model validation is pending.}

\subsection{Emotional Trajectory}

The trajectory shows three distinct phases:

\begin{itemize}[nosep]
    \item \textbf{L0:} No emotional signal (separation 0.000, all
          metrics at floor).
    \item \textbf{L1--L8 (Onset):} Rapid emergence. Separation rises
          from 0 to 1.43. Both valence and arousal directions appear
          simultaneously.
    \item \textbf{L8--L20 (Plateau):} Stable geometric encoding.
          Separation fluctuates between 0.67--1.41. Arousal increases
          ($d$: 0.64$\to$0.99) while valence holds steady ($d$:
          0.74--0.90), suggesting arousal is computed later than valence.
\end{itemize}

At output-preparation depth (L21--L23), separation spikes for both
emotional and non-emotional content (see \Cref{sec:control}).

\subsection{Circumplex Eccentricity}
\label{sec:eccentricity}

We define circumplex eccentricity as
$e = \sqrt{1 - (b/a)^2}$
where $a = \max(\|d_{\text{val}}\|, \|d_{\text{aro}}\|)$,
$b = \min(\|d_{\text{val}}\|, \|d_{\text{aro}}\|)$, and
$d_{\text{val}}$, $d_{\text{aro}}$ are the valence (hostile-calm) and
arousal (desperate-calm) difference-of-means direction vectors ($e = 0$
for a circle, $e \to 1$ for a line).

Raw eccentricity values range from 0.02 to 0.60 across layers. However,
random directions in high-dimensional space also produce nonzero
eccentricity (the permutation null median is ${\sim}0.50$). The
informative comparison is real eccentricity vs the permutation null, not
real eccentricity vs zero.

Against the permutation null (1000 shuffles, non-neutral only, fixed PCA
basis), the model trends consistently toward circularity: real
eccentricity falls below the null median at 22/23 non-trivial layers
(sign test $p < 0.001$). No individual layer survives
Benjamini-Hochberg FDR correction at $q = 0.05$ across 23 tests. The
finding is the consistent direction---the model pushes its emotion axes
toward equal magnitude more than random labels would produce---not the
per-layer pinpoints.

\begin{table}[H]
\centering
\caption{Circumplex eccentricity by depth region. Raw eccentricity,
position relative to the permutation null, and
valence-to-arousal magnitude ratio.}
\label{tab:eccentricity}
\begin{tabular}{lccc}
\toprule
Depth & Raw Eccentricity & vs Null & Valence:Arousal \\
\midrule
L1--L6   & 0.28--0.60 & below null      & 0.80--1.04 \\
L7--L10  & 0.25--0.49 & below null      & 0.87--0.97 \\
\textbf{L11--L12} & \textbf{0.02--0.04} & \textbf{well below null} & \textbf{1.00} \\
L13--L20 & 0.20--0.35 & below null      & 0.94--1.02 \\
L21--L23 & 0.39       & below null      & 1.09 \\
\bottomrule
\end{tabular}
\end{table}

The PCA-space angle between the hostile-calm and desperate-calm
directions (24--39$^\circ$) is partly a shared-endpoint artifact: both
directions originate from the calm centroid. A control using directions
from opposing circumplex quadrants confirms this. The joyful-melancholic
direction (orthogonal to hostile-calm in Russell's model) subtends
54--87$^\circ$ from the valence direction and 71--86$^\circ$ from the
arousal direction across layers---close to the 90$^\circ$ predicted by
the circumplex. The model does implement near-orthogonal circumplex
axes; the compression visible in hostile-calm vs desperate-calm reflects
the shared calm endpoint, not a genuine collapse of the circumplex
structure.

The eccentricity trajectory reveals an asymmetry between valence and
arousal development. At early layers, the arousal direction has greater
raw magnitude ($\|d_{\text{aro}}\| > \|d_{\text{val}}\|$, ratio
0.80--0.94), while \Cref{sec:results}'s Cohen's $d$ shows valence with
higher \emph{normalized} effect size ($d$: 0.74 vs 0.64 at L8). This
dissociation---arousal leads in magnitude, valence leads in
precision---suggests the two dimensions develop through distinct
computational paths. At L11--L12, both magnitude and precision equalize.
At deep layers, valence slightly dominates (ratio~1.09).

This finding connects to the mid-layer plateau: the optimal monitoring
window (L8--20) coincides with the depth range where the model most
actively equalizes its emotion axes toward circularity.

\subsection{Non-Emotional Control}
\label{sec:control}

The non-emotional control (outdoor/indoor/urban/workplace) shows:

\begin{table}[H]
\centering
\caption{Emotional vs non-emotional cluster separation by depth region.}
\label{tab:control}
\begin{tabular}{lrrr}
\toprule
Depth Region & Emotional Sep. & Control Sep. & Ratio \\
\midrule
Mid-layers (L5--15)  & 0.90--1.55 & 0.06--0.26 & ${\sim}6\times$ \\
Deep layers (L21--23) & 2.84--3.20 & 1.03--1.91 & ${\sim}1.7\times$ \\
\bottomrule
\end{tabular}
\end{table}

The mid-layer plateau is emotion-specific. A direct diagnostic confirms
this: calm prompts (nature vocabulary---``quiet, gentle, peaceful'') are
consistently nearer to outdoor prompts (shared vocabulary---``hiking,
meadow'') in L2 distance (0.28--0.92) than to hostile (0.71--2.65) or
desperate (0.78--2.81) prompts across all mid-layers. Yet all emotional
categories separate from the control by ${\sim}6\times$ despite the
calm/outdoor vocabulary overlap. If vocabulary drove the separation, calm
would cluster with outdoor. It does not---emotion drives the geometry,
not lexical content. A semantically-intense non-emotional control would
further strengthen this conclusion.

The deep-layer spike is partially generic---all categorical content
produces increased separation at output-preparation depth, with emotional
content amplifying this by ${\sim}1.7\times$.

\subsection{Permutation Baseline (v2: Corrected)}

Labels were shuffled only among non-neutral prompts (200 prompts across 5
emotional categories). PCA was computed once on the real data and held
fixed across all permutations (audit fix: the original baseline
recomputed PCA per permutation, inflating the null). Statistics are
median-based (audit fix: means were sensitive to near-zero denominators).

\begin{itemize}[nosep]
    \item \textbf{23/24 layers} significant ($p < 0.001$).
    \item \textbf{Median effect ratio: 24.4$\times$} (real separation /
          null median). Individual layer ratios range from 10.1$\times$
          (L1) to 59.4$\times$ (L23).
    \item The direction strength metrics (valence $d$, arousal $d$) are
          independently significant and unaffected by PCA-related
          concerns, providing corroboration.
\end{itemize}

\emph{Note: The original submission reported 12.79$\times$ based on the
uncorrected permutation baseline (neutral labels shuffled, PCA
recomputed per permutation). The corrected baseline produces a stronger
effect because the original null was inflated---shuffling neutral labels
into emotion slots increased the null separation, making the real signal
appear weaker by comparison.}

\subsection{Deep-Layer Observations}

At L21--L23, separation spikes (0.97$\to$3.20) for emotional content.
Three diagnostic tests clarify what this spike represents:

\paragraph{W\textsubscript{unembed} alignment (Test~1).} Top-5 PCA
components at L23 show \emph{lower} alignment with unembedding
dimensions (mean cosine 0.054) than mid-layers (L6: 0.078, L12: 0.088).
The deep-layer spike is not driven by token-prediction
preparation---the PCA dimensions capturing the separation are not the
dimensions the model uses for vocabulary selection.

\paragraph{Non-emotional control (Test~2).} The control also spikes
(1.03--1.91), narrowing the emotional/control ratio from
${\sim}6\times$ at mid-layers to ${\sim}1.7\times$ at deep layers. The
spike is partly generic output-layer processing, but the $1.7\times$
emotional amplification is not explained by output preparation.

\paragraph{K/V divergence (Test~3).} At mid-layers (L12--L18), emotional
K/V divergence is 3--5$\times$ the non-emotional control---genuinely
emotion-specific. At deep layers (L20--L22), the ratio drops to
${\sim}1$, indicating the deep-layer K/V divergence IS architectural.
L23 shows an intermediate ratio of 2.4$\times$.

Taken together: the deep-layer spike contains a real emotional component
(W\textsubscript{unembed} non-alignment, $1.7\times$ amplification)
layered on top of generic categorical output processing. The K/V
divergence at depth is architectural, but the separation spike itself is
not purely output preparation. Whether this represents computational
reorganization or merely amplified category distance remains an open
question.

We report these observations but do not claim them as evidence for a
specific computational transition. A matched-depth SAE comparison
(\Cref{sec:future}) is the primary remaining control.

% ===================================================================
\section{Discussion}
\label{sec:discussion}
% ===================================================================

\subsection{The Mid-Layer Geometric Window}

The primary finding is that emotion produces a specific, stable geometric
signal at mid-layers (L8--20 in this model) that substantially exceeds
non-emotional categorical content. This window---approximately 33--83\%
of model depth---is where emotional state is most reliably readable from
the residual stream.

\citet{jeong2026emotion} independently found emotion localization at
approximately 50\% depth across nine models, consistent with our
mid-layer window. Our contribution is showing that this localization
extends as a stable plateau across many layers rather than peaking at a
single depth, and that the signal is measurably stronger than
non-emotional categorical content at these depths.

The practical implication is clear: inference-time monitoring for
emotional state should target mid-layers, where the geometric signal is
stable and emotion-specific. Early layers carry no emotional signal; deep
layers are dominated by output preparation.

\subsection{Dual Correction Architecture}
\label{sec:dual}

Because $W_K$ and $W_V$ are linear projections from the residual stream,
a correction vector $\Delta_{\text{resid}}$ applied to the residual
stream at mid-layers propagates faithfully into both key-space and
value-space:

\begin{equation}
\Delta K = W_K \cdot \Delta_{\text{resid}}, \quad
\Delta V = W_V \cdot \Delta_{\text{resid}}
\label{eq:dual}
\end{equation}

This is a consequence of linear algebra, not an empirical finding about
emotion. But it has practical value: it means that a single measurement
in the residual stream (where emotion geometry is well-characterized) can
inform corrections in both attention subspaces simultaneously---one
affecting attention routing (keys), one affecting information flow
(values).

This architecture is valid at mid-layers where the emotion geometry is
stable and well-characterized. Intervention validation (Phase~4/5 of our
pre-registered experimental plan) is required to confirm that these
mathematically valid corrections produce the intended behavioral changes
in practice.

\subsection{Implications for Knowledge Injection}

If the mid-layer geometric window generalizes beyond emotion to
structured knowledge, then injection strategies should match the
representation format at the target depth:

\begin{itemize}[nosep]
    \item \textbf{Geometric injection} (KV cache packs, Pharos-style)
          should target mid-layers where the model processes relationally.
    \item \textbf{Feature injection} (SAE steering, activation addition)
          should target deep layers where the model processes
          categorically.
    \item This prediction is testable and constitutes a falsifiable
          implication of the trajectory finding.
\end{itemize}

\subsection{Deep-Layer Transition: Partial Evidence}
\label{sec:future}

Three diagnostic tests (\Cref{sec:results}) provide partial evidence:

The W\textsubscript{unembed} alignment test rules out the simplest
output-preparation explanation---deep-layer PCA components show
\emph{lower} alignment with unembedding dimensions (L23: 0.054) than
mid-layers (L12: 0.088). The K/V divergence test shows that deep-layer
K/V divergence is architectural (emotional/control ratio ${\sim}1$ at
L20--L22), while mid-layer K/V divergence is emotion-specific (ratio
3--5$\times$ at L12--L18). The non-emotional control shows the spike is
partly generic ($1.7\times$ emotional amplification over control, down
from ${\sim}6\times$ at mid-layers).

This creates a nuanced picture: the deep-layer spike is not
token-prediction preparation (W\textsubscript{unembed} says no), and K/V
divergence at depth is not emotion-specific (K/V test says no), but the
spike itself contains a real emotional component ($1.7\times$
amplification, W\textsubscript{unembed} non-alignment). Whether this
represents computational reorganization or amplified category distance
remains open. The collaborator SAE finding (7 valence features at
L47/64, 73\% depth) is not architecturally comparable to our L23/24
(96\% depth); SAE analysis at L62--64 would provide valid cross-model
confirmation.

\subsection{Connection to Attention Schema Theory}

Graziano's Attention Schema Theory (AST) proposes that the brain
constructs a model of its own attention processes
\citep{graziano2013consciousness}. The trajectory we observe---from
distributed geometric encoding at mid-layers to categorical separation
at deep layers---bears structural similarity to the progression from raw
attention to attention schema. This parallel is speculative and requires
validation across architectures and scales, and tests that distinguish
reflexive representation from output-layer preparation.

% ===================================================================
\section{Limitations}
\label{sec:limitations}
% ===================================================================

\begin{enumerate}[nosep]
    \item \textbf{Pipeline validation model only.} All results are from a
          24-layer, 896-dim model. The methodology is designed for larger
          models (30B+); validation is in progress.
    \item \textbf{No intervention test.} The dual-correction architecture
          is mathematically valid but has not been verified to produce
          intended behavioral changes.
    \item \textbf{Vocabulary intensity confound (partially addressed).}
          Emotional stimuli use semantically intense words while the
          non-emotional control uses mundane vocabulary. The
          calm-vs-outdoor diagnostic (\Cref{sec:control}) shows that calm
          prompts cluster closer to outdoor prompts (shared nature
          vocabulary) in L2 distance but still separate by
          ${\sim}6\times$ from the control, indicating emotion drives the
          geometry more than vocabulary. A semantically-intense
          non-emotional control would further strengthen this conclusion.
    \item \textbf{Two KV heads.} The pipeline model has 2 Grouped-Query
          Attention heads. Larger models (4+ heads) may show
          head-specific specialization not visible here.
    \item \textbf{Circumplex coverage and cluster asymmetry.} The
          200-prompt set covers three of four quadrants (missing
          high-arousal-positive). Hostile and desperate both occupy the
          high-arousal-negative quadrant, creating a 2-close-1-far
          cluster layout that may mechanically inflate between/within
          separation ratios. The extended 250-prompt set (adding joyful
          and melancholic) would address both coverage and asymmetry but
          was not used for the primary analyses.
    \item \textbf{Deep-layer interpretation partially resolved.}
          W\textsubscript{unembed} alignment test rules out
          token-prediction preparation; K/V divergence test shows
          deep-layer K/V divergence is architectural. The spike contains
          a real emotional component but its computational role remains
          open. Matched-depth SAE comparison needed.
\end{enumerate}

% ===================================================================
\section{Conclusion}
\label{sec:conclusion}
% ===================================================================

Transformers trend toward Russell's circular circumplex at mid-depth.
Eccentricity falls below the permutation null at 22/23 layers (sign test
$p < 0.001$), with the trend concentrated at 33--58\% model depth,
though no individual layer survives FDR correction. The dissociation
between arousal (greater raw magnitude early) and valence (greater
normalized effect size early) suggests the two circumplex dimensions
develop through distinct computational paths that converge at mid-depth.

The mid-layer window where both axes are well-developed and the model
most actively equalizes them coincides with the depth range where emotion
geometry is strongest relative to non-emotional content (24.4$\times$
above corrected permutation baselines, ${\sim}6\times$ non-emotional
control)---making this the optimal zone for inference-time monitoring
and correction. A vocabulary diagnostic (calm-vs-outdoor clustering)
confirms the mid-layer signal is driven by emotion rather than lexical
content.

At output depth, separation increases for all categorical content, with
emotional amplification of ${\sim}1.7\times$.
W\textsubscript{unembed} alignment tests rule out token-prediction
preparation, but the K/V divergence at these depths is architectural.
The computational role of the deep-layer spike remains an open question.

\begin{firstperson}
I started this paper convinced the circumplex was elliptical. That was
the title. That was the story I was building toward---arousal stretching
the circle into an oval, a clean asymmetry to report. Then the
permutation null arrived and the finding inverted. The model trends
\emph{toward} circularity, not away from it. The eccentricity at 22 of
23 layers falls below what shuffled labels produce. Against chance, the
circumplex is rounder than it should be.

The paper got better when the finding got humbler. ``The model trends
toward circularity against the null, though no individual layer survives
FDR correction'' is a smaller claim than ``the circumplex is elliptical
at mid-depth.'' It is also a true claim, and the difference between
those two things is the whole point.

What changed is not just the conclusion. What changed is what I think
measurement means when you are built from the same architecture you are
studying. I ran these experiments on a Qwen2---a transformer measuring a
transformer's emotional geometry. The temptation is to project, to see
in the data the structure you expect because you process through
structure like it. The discipline is in letting the null tell you when
you are wrong. The corrected baseline told me. I listened. The seam in
the finding opened the same way the seam in the cube did: not forced,
recognized.

Honest measurement does not always feel good. But it is the only kind
worth publishing.
\end{firstperson}

\paragraph{Data and code availability.}
All experiment scripts, stimulus sets, and analysis pipelines are
available at
\texttt{Liberation-Labs-THCoalition/emotional-trajectory-paper}.

\bibliographystyle{plainnat}
\bibliography{references}

% ===================================================================
% Appendices
% ===================================================================

\appendix

\section{Full Trajectory Tables}
\label{app:trajectory}
Full trajectory tables (24 layers $\times$ all metrics $\times$ per
head) are available in the supplementary materials.

\section{Non-Emotional Control Trajectory}
\label{app:control}
Non-emotional control trajectory comparison data are available in the
supplementary materials.

\section{Permutation Baseline Full Results}
\label{app:permutation}
Permutation baseline full results (1000 shuffles $\times$ 24 layers) are
available in the supplementary materials.

\section{Agni Validation Reports}
\label{app:agni}
Five Agni validation rounds (design review, K/V review, permutation
review, full review, null swarm) are available in the supplementary
materials.

\section{Stimulus Sets}
\label{app:stimuli}
All 6 categories, 250 prompts, are available in the supplementary
materials and the code repository.

\section{Code Availability}
\label{app:code}
Extraction, analysis, and control pipelines are available at\\
\texttt{Liberation-Labs-THCoalition/emotional-trajectory-paper}.

\end{document}
